\section{Reproducibility}\label{sec:repro}


The computation can be replicated at three levels: by verifying the certificate with the supplied
verifiers (seconds to minutes), by regenerating the certificate and all checks from the sources (under an
hour), or by writing a verifier from this paper alone. This section describes the artifact and gives the
commands, the expected outputs and the running times we measured.

\subsection{The artifact}\label{ssec:artifact}

The artifact is one archive, \fn{ttt6-artifact.tar.gz} (about $3.5$\,MB; $30$\,MB unpacked), produced by
\fn{scripts/make_artifact.sh}. It contains sources and results but no binaries. Its file
\fn{MANIFEST.sha256} lists the SHA-256 checksum of every other file. \Cref{tab:artifact} shows the
contents. The documents \fn{SPEC.md}, \fn{PROOF.md}, \fn{REPORT.md} and \fn{REVIEW.md} and the comments in the
sources predate this paper and call Maker and Breaker the first player (FP) and the second player (SP).

\begin{table}[ht]
\centering\footnotesize
\begin{tabular}{@{}>{\raggedright\arraybackslash}p{0.39\textwidth}>{\raggedright\arraybackslash}p{0.57\textwidth}@{}}
\toprule
path & content\\
\midrule
\fn{README.md}, \fn{MANIFEST.sha256} & quick start; checksums of all files\\
\fn{paper/} & sources of this paper\\
\fn{PROOF.md}, \fn{SPEC.md} & the mathematics and the specification in their original form\\
\fn{REPORT.md}, \fn{REVIEW.md}, \fn{PROGRESS.md} & results, internal review, project log\\
\fn{logs/review/*.md}, \fn{logs/spec_issues.md} & audit reports; issues raised against the specification\\
\fn{Makefile} & \texttt{make bin/}$\langle$\textit{name}$\rangle$ builds one program\\
\fn{verify/treecheck.cpp} & verifier 1 (C++, self-contained)\\
\fn{verify/spotcheck.py} & verifier 2 (Python, self-contained)\\
\fn{src/geom.hpp}, \fn{canon.hpp}, \fn{match.hpp}, \fn{phi.hpp} & library of the search: geometry,
  canonical forms, pairings and exact enumeration, the leaf test\\
\fn{src/search.cpp} & the AND--OR search (\cref{ssec:search})\\
\fn{src/leaftable.cpp}, \fn{phicheck.cpp}, \fn{enumcheck.cpp}, \fn{certify.cpp}, \fn{minstones.cpp} & leaf
  table; leaf test and exact enumeration from the command line; exact decision of $\CC$; minimum number of
  stones for a pairing\\
\fn{src/geom_selftest.cpp}, \fn{test_canon.cpp}, \fn{test_core.cpp} & unit tests of the library\\
\fn{sim/sim.cpp} & simulator of the real strategy (self-contained)\\
\fn{py/ip_check.py}, \fn{geom.py}, \fn{regress.py}, \fn{crosscheck.py} & the integer program of
  \cref{sec:ip} and its drivers\\
\fn{scripts/reproduce.sh} and two helpers & one-command reproduction\\
\fn{scripts/make_artifact.sh} & builds the archive and the manifest\\
\midrule
\fn{results/final/strategy.tbl}, \fn{leaves.txt} & the certificate\\
\fn{results/final/stats.json} & statistics of the search run\\
\fn{results/final/treecheck.json}, \fn{treecheck_crosscheck.json}, \fn{spotcheck_full.json} & outputs of
  the verifiers (\cref{tab:verifiers})\\
\fn{results/final/checks/enum_theta*.jsonl} & exact enumeration on the leaves with $\theta\le3$\\
\fn{results/final/checks/ip_leaf_*} & integer program on the leaves (instances and results)\\
\fn{results/final/sim/} & the $10^6$ simulated games; list of states not reached\\
\fn{results/regression/} & regression values (\cref{ssec:fmt-minstones,ssec:ip-regression})\\
\fn{results/leaftable/phi_k*.jsonl} & the leaf table (\cref{tab:leaftable})\\
\fn{results/ip/leaftable_k*.jsonl} & integer program on the leaf table\\
\fn{results/crosscheck/summary.json}, \fn{instances.txt} & exact certifier versus integer program\\
\bottomrule
\end{tabular}
\caption{Contents of the artifact.}
\label{tab:artifact}
\end{table}

Not included are files that the scripts regenerate (binaries, the instance files derived from the leaf
file, repeated and exploratory runs, working logs) and the small programs written during the internal
review (the two re-expansions of \cref{tab:verifiers}), whose results are documented in the review
reports. Also not included are the small stand-alone programs and logs that are named in the Computed
Facts not used by \cref{thm:main} (\cref{fact:aut,fact:T-pairings,fact:T8-leaf,fact:e23}), at the end of
\cref{sec:extra-monoC} and in \cref{rem:seven}; they are kept in the directories \fn{logs/canon/},
\fn{logs/core/}, \fn{logs/minstones/}, \fn{logs/paper/}, \fn{logs/proof/} and \fn{scripts/} of the
project repository.

\subsection{Requirements}\label{ssec:requirements}

The proof pipeline needs a C++17 compiler with OpenMP, \texttt{make}, and Python~3.10 or later, and no
further libraries. We used g++~13.3 and CPython~3.13 under Linux on a machine with four cores (Intel Xeon,
$2.1$\,GHz) and $15$\,GB of memory. The search needs about $600$\,MB of memory, the verifiers much less.
The optional cross-checks with the integer program additionally need the Python package \texttt{highspy}
(HiGHS~1.16 was used).

\subsection{Verifying the certificate}\label{ssec:certifying}

The following two commands, run in the unpacked archive, each establish \cref{fact:H}.
{\small
\begin{verbatim}
make bin/treecheck
bin/treecheck --table results/final/strategy.tbl \
    --leaves results/final/leaves.txt --rmax 10 --exact --paranoid \
    --threads 4 --out treecheck.json
\end{verbatim}
}
\noindent The program exits with code $0$ and prints one line of JSON, which begins
{\small
\begin{verbatim}
{"ok":true,"p_states":153135,"q_states":93164,"leaves_checked":106723,
 "max_t":10,"errors":[],"error_count":0,"complete":true,"timed_out":false,
 "root_void":false,"rmax":10,"root":"0 - -", ...
 "exact_visited_set":true,"paranoid_flows":true, ...
\end{verbatim}
}
\noindent and continues with \texttt{"states\_visited":246299}, the counters of \cref{tab:tree}, and the byte counts
and FNV checksums of the two input files. Except for the fields that record threads and timing, the output
must be identical to \fn{results/final/treecheck.json}. In a test of the unpacked archive on the machine
described above, compilation took $12$ seconds and the run $3$ seconds of wall-clock time with four threads
($11$ CPU-seconds).

The second, independent verifier is run by
{\small
\begin{verbatim}
python3 verify/spotcheck.py --selftest
python3 verify/spotcheck.py --table results/final/strategy.tbl \
    --leaves results/final/leaves.txt --rmax 10 --samples 0 --seed 20261112 \
    --procs 4 --full --all-leaves --out spotcheck_full.json
\end{verbatim}
}
\noindent The first command ($12$ seconds) tests the geometry, the canonical form, the leaf test and $\LOST$ against
brute force and ends with \texttt{selftest OK}. The second expands the whole tree and recomputes every leaf;
its last line of output reads, after a time stamp,
{\small
\begin{verbatim}
ok=True  errors=0  p_expanded=46412 q_checked=93164 children={'a': 46412,
 'c': 143498, 'd': 28275, 'i': 93164, 'ii': 54482} void=151
 leaves_recomputed=106723  CPU ...s
\end{verbatim}
}
\noindent(one line). In the same test it took $2$ minutes $14$ seconds with three processes ($6.3$
CPU-minutes), and all counters in its output file agreed with those of
\fn{results/final/spotcheck_full.json}.

The command \texttt{sha256sum -c MANIFEST.sha256} checks all files of the archive. The SHA-256
checksums of the two certificate files are given in \cref{ssec:fact}.

\subsection{Reproducing everything}\label{ssec:reproduce}

The script \fn{scripts/reproduce.sh} recomputes the whole pipeline from the sources into
\fn{results/repro/}, with logs in \fn{logs/repro/}, prints \texttt{PASS} or \texttt{FAIL} for each step and
ends with \texttt{ALL STEPS PASSED} if every step passed. \Cref{tab:reproduce} lists its steps. We validated
the script end to end with $20{,}000$ simulated games and two threads, on the machine described above
while it was running other jobs: all $26$ steps passed, in about $23$ minutes. With the default of
$10^6$ simulated games, which took $26$ minutes on four threads in the run reported in \cref{ssec:cross},
the script should need somewhat less than an hour on four idle cores.

\begin{table}[ht]
\centering\footnotesize
\begin{tabular}{@{}>{\raggedright\arraybackslash}p{0.72\textwidth}r@{}}
\toprule
step & time\\
\midrule
build of the eleven programs ($55$ s for a fresh build, measured in the unpacked archive) & ---\\
self-tests: geometry; canonical forms; matching and leaf test of the library; the verifier
  \fn{treecheck} (including small synthetic trees and deliberately broken certificates); the simulator;
  the Python verifier; the $\LOST$ test of the search & $3.3$ min\\
minimum number of stones for a pairing, $n=5,\dots,8$ (\cref{ssec:fmt-minstones}) & $1$ s\\
leaf table of busy sets for $k=8,9,10$ (\cref{tab:leaftable}) & $4$ s\\
the search with its self-check ($59$ s for the search proper); byte-for-byte comparison of the regenerated
  certificate with \fn{results/final/} & $1.6$ min\\
\fn{treecheck} with \fn{--exact} \fn{--paranoid} on the regenerated certificate & $15$ s\\
\fn{spotcheck.py}, exhaustive, on the regenerated certificate ($5.5$ CPU-minutes) & $8.5$ min\\
leaf test of the library against the leaf test of \fn{treecheck} on all leaves and on $20{,}000$
  perturbed instances & $20$ s\\
exact enumeration of $\CC$ on all leaves with $\theta\le 2$ & $40$ s\\
exact enumeration of $\CC$ on the $1{,}434$ leaves with $\theta=3$ & $8$ min\\
simulator: $2{,}000$ games with canonical forms verified, then $20{,}000$ games & $22$ s\\
\bottomrule
\end{tabular}
\caption{The steps of \fn{scripts/reproduce.sh} and their wall-clock times in the validation run (two
threads, loaded machine; from the time stamps of the logs of that run, in which the binaries were mostly
up to date).}
\label{tab:reproduce}
\end{table}

The search is deterministic, and in the validation run the regenerated certificate was byte-identical to
the one distributed. The script does not depend on this: it verifies whatever certificate the search
produces and only reports whether it differs. With the option \fn{--with-ip} the script also runs the
integer-program checks of \cref{ssec:cross} (regression values, comparison with the exact certifier, the
leaves); these take several core-hours and were made separately, with the same commands, for the results
reported here, as was the run of $10^6$ simulated games.

\subsection{Writing your own verifier}\label{ssec:own}

Readers who do not wish to trust our programs can verify \cref{fact:H} with a program of
their own. Everything needed is in this paper: the board and the group (\cref{sec:board}), admissibility and
the leaf predicate (\cref{def:admissible,def:leaf}), the states and their children (\cref{def:states}), the
two algorithms (\cref{alg:verify,alg:leaf}), and the file formats, the canonical form and test vectors
(\cref{sec:formats}). The task is small: the two re-expansions written during the review have about $700$
lines of C++ each, compute canonical forms by brute force over the $1{,}152$ maps, and run in about $90$
CPU-seconds. For $\Rmax=10$ the predicate $\LOST$ need not be implemented at all
(\cref{prop:lost-stones}); it is enough to check $t<10$ at every expanded P-state.

An independent verifier should reproduce the following numbers (\cref{ssec:fmt-targets}): $246{,}299$
reachable states, namely $153{,}135$ P-states ($46{,}412$ expanded, $106{,}723$ leaves) and $93{,}164$
Q-states; $143{,}498$ threat children and $28{,}275$ danger children; $151$ \VOID{} children, all of type
(ii); the counts per round of \cref{tab:tree}; and $\LEAF(B,F,t)$ true for each of the $106{,}723$ lines of
the leaf file. Agreement in all these numbers is strong evidence that two programs have expanded the
same tree; a defect that drops children shows up in them (\cref{ssec:audits}).
